\newpage
\chapter{Functional Analysis}

\section{Contractions and Banach}

\subsection{Local Convergence}

Given \(x_k \in X \subseteq \R^n\), we define \emph{local convergence} for the iteration \(x_{k + 1} = \Phi(x_k)\) as 

\[
    \lim_{k \to \infty} x_k = x \ \forall x_0 \in X.
\]

\subsection{Contraction}

\(\Phi\) is considered a contraction according to \(\metric{\cdot}\) on \(X \subseteq \R^n\) if 

\[
    \Phi: X \to X
\]

and there is an \(\alpha < 1\) independent from \(x,y\) such that 

\[
    \metric{\Phi(x) - \Phi(y)} \le \alpha \metric{x - y} \ \forall x,y \in X
\]

To check if a function is a contraction, on a certain set or interval, we need to check the following properties:

\begin{itemize}

    \item Check that \(f:X \to X\), i.e. that the function at its local extremes and endpoints remains in the interval. This 
    accomplished by the derivative test or if the function is monotone, by checking the function at the endpoints of the interval.

    \item Estimating \(\alpha\) using the definition of the contraction and using approximations \(\ge\). Alternatively, we can use the \emph{Mean Value Theorem} with 
    \(\max_{\xi \in X} \metric{f'(x)} \le \alpha\) via a derivative test using \(\xi\) as our variable.

\end{itemize}

\textbf{Example I:}

Check if \(f(x) = \sqrt{x}\) is a contraction on \([1,\infty)\) with the euclidean metric. 

Because of \(f(x)\) being monotone, we just need to check the function at the endpoints of the interval.

\[
    f(1) = 1 \quad \text{and} \quad f(\infty) = \infty
\]

Thus, our first property is satisfied. Now we check the second property:

\begin{align*}
    \metric{\sqrt{x} - \sqrt{y}} &= \frac{\metric{x - y}}{\metric{\sqrt{x} + \sqrt{y}}} \\ 
    &\le \frac{\metric{x - y}}{\metric{1 + 1}} = \frac{1}{2} \metric{x - y}
\end{align*}

Hence \(\alpha = \frac{1}{2}\).

\textbf{Example II:}

\(\Phi(x) = x^3 -3x^2 + 3x\) on \(X = [\frac{1}{2}, \frac{3}{2}]\).

We take the first derivatives: 

\[
    \Phi'(x) = 3x^2 - 6x + 3 \quad \text{and} \quad \Phi''(x) = 6x - 6
\]

Then we perform our tests:

\(\Phi'(x) = 0\):

\[
    \Phi'(x) = 0 \iff 3x^2 - 6x + 3 = 0 \iff x = 1 > 0
\]

\(\Phi''(x) \ne 0\):

\[
    \Phi''(1) = 6 - 6 = 0
\]

We do not have any local extremes, hence we check the function at the endpoints of the interval:

\[
    \Phi\left(\frac{1}{2}\right) = \frac{11}{8} \in X
\]

\[
    \Phi\left(\frac{3}{2}\right) = \frac{11}{8} \in X
\]

Hence \(\Phi(X) \subseteq X\). Now we check the second property:

\[
    \alpha \le \max_{x \in X} \metric{\Phi'(x)} 
\]

\subsection{The Fix Point Theorem of Banach}

Given a closed set \(X \subseteq \R^n\) with the contraction \(\Phi : X \to X\) and 

\[
    \metric{\Phi(x) - \Phi(y)} \le \alpha \metric{x - y} \ \forall x,y \in X
\]

such that \(\alpha < 1\) is independent. Then \(\Phi\) has exactly one fix-point \(x = \Phi(x)\) in \(X\). Our
\emph{Picard iteration} converges for all \(x_0 \in X\) towards \(x\). Also the following approximations hold:

\[
    \metric{x - x_k} \le \frac{\alpha^k}{1 - \alpha} \metric{x_1 - x_0} \quad \text{a-priori}
\]

\[
    \metric{x - x_k} \le \frac{\alpha}{1 - \alpha} \metric{x_k - x_{k - 1}} \quad \text{a-posteriori}
\]

\textbf{Proof:}

Let \(X\) be a complete space and \(\Phi\) be a contraction: 

\[
    \metric{\Phi(x) - \Phi(y)} < \alpha \metric{x - y}
\]

We will show that \(x_k\) is a Cauchy sequence:

\[
    \metric{x_{k + 1} - x_k} = \metric{\Phi(x_k) - \Phi(x_{k - 1})} \le \alpha \metric{x_k - x_{k - 1}}
\]

Because of the contraction property and our recursive definition, we can find the following:

\begin{align*}
    \metric{f(x_2) - f(x_1)} &< \alpha\metric{x_2 - x_1} \\ 
    \metric{x_3 - x_2} &< \alpha\metric{x_2 - x_1} \\ 
    \metric{x_4 - x_3} &< \alpha \metric{x_3 - x_2} 
\end{align*}

Using the previous inequality we get

\[
    \metric{x_3 - x_2} < \alpha\metric{x_2 - x_1} \iff \alpha \metric{x_3 - x_2} < \alpha^2\metric{x_2 - x_1} 
\]

Now by applying it multiple times we get 

\begin{align*} 
    \metric{x_4 - x_3} &< \alpha \metric{x_3 - x_2} < \alpha^2 \metric{x_2 - x_1}\\ 
    &\dots \\ 
    \metric{f(x_{k + 1}) - f(x_{k})} &< \alpha^k \metric{x_2 - x_1} \\ 
\end{align*}

We use the following approximations for \(n > m\):

\begin{align*}
   \metric{x_n - x_m} = \metric{x_n + x_{n - 1} - x_{n - 1} + \dots + x_{m - 1} - x_{m - 1} - x_m} &\le \metric{x_n - x_{n - 1}} + \dots + \metric{x_{m + 1} - x_m} \\
    &\le \alpha^{n - 1} \metric{x_2 - x_1} + \dots + \alpha^m \metric{x_2 - x_1} \\
    &= (\alpha^{n - 1} + \dots + \alpha^m) \metric{x_2 - x_1} \\
    &= \alpha^m \sum_{i = 0}^{n - m - 1} \alpha^i \metric{x_2 - x_1} \\
    &= \alpha^m \frac{1 - \alpha^{n - m}}{1 - \alpha} \metric{x_2 - x_1} \\ 
    &< \frac{\alpha^m}{1 - \alpha} \metric{x_2 - x_1}
\end{align*}

Hence, our sequence is Cauchy and because \(X\) is complete, we know that \(\lim_{k \to \infty} x_k = p\) for some \(p \in X\).

We take the limit

\[
    \lim_{k \to \infty} \metric{f(x_{k + 1}) - f(x_{k})} < \alpha^k \metric{x_2 - x_1} = 0
\]

Because \(X\) is complete we know that \(\lim_{k \to \infty} f(x_k) = p\), hence \(f(p) = p\)

\QED

\subsection{Linear Convergence}

For the iteration \(\Phi(x_k) = x_{k + 1}\), we define \emph{linear convergence} as 

\[
    \metric{e_{k + 1}} \le c \metric{e_k} \ \forall k, \ 0 \le c < 1.
\]

The iteration is linearly convergent if such constant \(c\) exists.

\subsection{Convergence With Order \texorpdfstring{m}{}}

\emph{Convergence With Order \(m \ge 1\)} is given for an iteration \(\Phi(x_k) = x_{k + 1}\) if, \(\Phi(x)\) is 
linearly convergent and \(0 \le c\) is independent from \(k\) and 

\[
    \lim_{k \to \infty} \metric{e_K} = 0 \quad \text{and} \quad \metric{e_{k + 1}} \le c \metric{e_K}^m \ \forall k
\]

\subsection{Convergence Of The Differentiable Iteration}

If given an \(m \ge 2\), a fix-point \(x\) of \(\Phi: \R \to \R\) such that \(\Phi\) is \(m\)-times differentiable on \(x\) with the 
property 

\[
    \Phi'(x) = \Phi''(x) = \dots = \Phi^{(m - 1)}(x) = 0,
\]

then there exists a closed interval \(X = [x - \delta, x + \delta], \delta > 0\) such that the iteration 

\[
    x_{k + 1} = \Phi(x_k)
\]

for all \(x_0 \in X\) with at least order \(m\) converges.

This means that if we want to check the convergence of an iteration we can check the derivatives of the function at the fix-point and if they are zero up to a certain order, then 
we can conclude that the iteration converges with at least that order.




